%%%%%%%%%%%%%%%%%%%%%%%%%%%%%%%%%%%%%%%%%%%%%%%%%%%%%%%%%%%%%%%%%%%%%%%%%%%%%%%%
\begin{frame}[fragile]\frametitle{}
\begin{center}
{\Large Subgraphs and Multi-Agent Systems}
\end{center}
\end{frame}
%%%%%%%%%%%%%%%%%%%%%%%%%%%%%%%%%%%%%%%%%%%%%%%%%%%%%%%%%%%%%%%%%%%%%%%%%%%%%%%%
\begin{frame}[fragile]\frametitle{Advanced Concept: Subgraphs and Modularity}
      \begin{itemize}
        \item \textbf{What}: Embed complete graphs as nodes within parent graphs
        \item \textbf{Benefits}:
        \begin{itemize}
            \item Modular, reusable workflow components
            \item Hierarchical organization of complex systems
            \item Encapsulation and separation of concerns
        \end{itemize}
        \item \textbf{Example Hierarchy}:
        \begin{itemize}
            \item Parent: Customer service orchestrator
            \item Subgraph 1: Email classification workflow
            \item Subgraph 2: Ticket routing workflow
            \item Subgraph 3: Response generation workflow
        \end{itemize}
        \item Each subgraph is self-contained with own state and logic
        \item Parent graph coordinates between subgraphs
        \item Enables building complex systems from simple components
      \end{itemize}
\end{frame}
%%%%%%%%%%%%%%%%%%%%%%%%%%%%%%%%%%%%%%%%%%%%%%%%%%%%%%%%%%%%%%%%%%%%%%%%%%%%%%%%
\begin{frame}[fragile]\frametitle{Subgraph Architecture}
A subgraph is a compiled graph that is used as a node in a larger graph.

      \begin{itemize}
        \item Break complex workflows into manageable components
        \item Modular design enhances code reusability
        \item Independent subgraphs improve maintainability
        \item Easy testing and debugging of individual components
        \item Scalable architecture for adding new features
        \item Support composition and interaction between subgraphs
        \item Enable hierarchical workflow organization
      \end{itemize}
\end{frame}
%%%%%%%%%%%%%%%%%%%%%%%%%%%%%%%%%%%%%%%%%%%%%%%%%%%%%%%%%%%%%%%%%%%%%%%%%%%%%%%%
\begin{frame}[fragile]\frametitle{Subgraph Implementation: A Basic Subgraph}

A subgraph is built exactly like any graph: a \texttt{StateGraph}, nodes, edges, then \texttt{compile()}.

% [code-font: forced \scriptsize | overflow at the default size]
\begin{lstlisting}[language=Python, basicstyle=\scriptsize\ttfamily]
from typing_extensions import TypedDict
from langgraph.graph import StateGraph, START, END

class ContentState(TypedDict):
    topic: str
    draft: str
    final_content: str

def generate_content(state: ContentState):
    return {"draft": llm.invoke(f"Write a short post on {state['topic']}").content}

def review_content(state: ContentState):
    return {"final_content": llm.invoke(f"Improve this post:\n{state['draft']}").content}

content_builder = StateGraph(ContentState)
content_builder.add_node("generate", generate_content)
content_builder.add_node("review", review_content)
content_builder.add_edge(START, "generate")
content_builder.add_edge("generate", "review")
content_builder.add_edge("review", END)

content_graph = content_builder.compile()     # usable as a node
\end{lstlisting}
\end{frame}
%%%%%%%%%%%%%%%%%%%%%%%%%%%%%%%%%%%%%%%%%%%%%%%%%%%%%%%%%%%%%%%%%%%%%%%%%%%%%%%%
\begin{frame}[fragile]\frametitle{Subgraph Implementation: Own State Schema}

Each subgraph declares its own state; nodes return updates and never mutate the incoming state.

\begin{lstlisting}[language=Python]
class AnalyticsState(TypedDict):
    final_content: str
    metrics: dict

def process_analytics(state: AnalyticsState):
    # Return the changed field only: a new metrics dict
    return {"metrics": {"engagement": calculate_engagement(state["final_content"])}}

analytics_builder = StateGraph(AnalyticsState)
analytics_builder.add_node("analytics", process_analytics)
analytics_builder.add_edge(START, "analytics")
analytics_builder.add_edge("analytics", END)

analytics_graph = analytics_builder.compile()
\end{lstlisting}
\end{frame}
%%%%%%%%%%%%%%%%%%%%%%%%%%%%%%%%%%%%%%%%%%%%%%%%%%%%%%%%%%%%%%%%%%%%%%%%%%%%%%%%
\begin{frame}[fragile]\frametitle{Subgraph Implementation: Using Subgraphs in the Main Graph}

The parent shares state keys with its subgraphs (\texttt{final\_content} here), so compiled subgraphs plug in directly as nodes.

\begin{lstlisting}[language=Python]
class MarketingState(TypedDict):
    topic: str
    draft: str
    final_content: str
    metrics: dict

main_builder = StateGraph(MarketingState)

# Add the compiled subgraphs as nodes
main_builder.add_node("content", content_graph)
main_builder.add_node("analytics", analytics_graph)

# Connect them like any other nodes
main_builder.add_edge(START, "content")
main_builder.add_edge("content", "analytics")
main_builder.add_edge("analytics", END)

marketing_app = main_builder.compile()
result = marketing_app.invoke({"topic": "LangGraph"})
\end{lstlisting}
\end{frame}
%%%%%%%%%%%%%%%%%%%%%%%%%%%%%%%%%%%%%%%%%%%%%%%%%%%%%%%%%%%%%%%%%%%%%%%%%%%%%%%%
\begin{frame}[fragile]\frametitle{Subgraph Implementation: Data Passing Between Subgraphs}
      \begin{itemize}
        \item \textbf{Shared keys flow automatically}: whatever one subgraph writes to a shared key, the next one reads.
        \item \textbf{Different schemas}: call the subgraph from a wrapper node and map fields in and out.
        \item Keep the interface small: pass only the fields the next stage needs.
      \end{itemize}
\begin{lstlisting}[language=Python]
def prepare_data(state: MarketingState):
    # Reshape the output of one stage for the next stage
    return {"metrics": {"content_type": "post",
                        "length": len(state["final_content"])}}

def call_analytics(state: MarketingState):
    # Wrapper node: map parent fields to the child schema and back
    out = analytics_graph.invoke({"final_content": state["final_content"],
                                  "metrics": {}})
    return {"metrics": out["metrics"]}
\end{lstlisting}
\end{frame}
%%%%%%%%%%%%%%%%%%%%%%%%%%%%%%%%%%%%%%%%%%%%%%%%%%%%%%%%%%%%%%%%%%%%%%%%%%%%%%%%
\begin{frame}[fragile]\frametitle{Practical Case: Marketing Agent, Content Subgraph}

% [code-font: forced \scriptsize | overflow at the default size]
\begin{lstlisting}[language=Python, basicstyle=\scriptsize\ttfamily]
class CampaignState(TypedDict):
    audience: str
    platform: str
    goal: str
    generated_content: str
    final_content: str
    analytics: dict
    recommendations: str

def generate_content(state: CampaignState):
    prompt = (f"Target Audience: {state['audience']}\n"
              f"Platform: {state['platform']}\n"
              f"Campaign Goal: {state['goal']}\n"
              "Write the campaign post.")
    return {"generated_content": llm.invoke(prompt).content}

def optimize_content(state: CampaignState):
    # Adapt length, tone and hashtags to the platform
    optimized = optimize_for_platform(state["generated_content"], state["platform"])
    return {"final_content": optimized}

content_b = StateGraph(CampaignState)
content_b.add_node("generate", generate_content)
content_b.add_node("optimize", optimize_content)
content_b.add_edge(START, "generate")
content_b.add_edge("generate", "optimize")
content_b.add_edge("optimize", END)
content_sub = content_b.compile()
\end{lstlisting}
\end{frame}
%%%%%%%%%%%%%%%%%%%%%%%%%%%%%%%%%%%%%%%%%%%%%%%%%%%%%%%%%%%%%%%%%%%%%%%%%%%%%%%%
\begin{frame}[fragile]\frametitle{Practical Case: Marketing Agent, Analytics Subgraph}

\begin{lstlisting}[language=Python]
def analyze_performance(state: CampaignState):
    m = calculate_metrics(state["final_content"])
    return {"analytics": {"engagement_score": m["engagement"],
                          "reach_prediction": m["reach"],
                          "conversion_estimate": m["conversion"]}}

def generate_recommendations(state: CampaignState):
    recs = generate_improvements(state["analytics"], state["goal"])
    return {"recommendations": recs}

analytics_b = StateGraph(CampaignState)
analytics_b.add_node("analyze", analyze_performance)
analytics_b.add_node("recommend", generate_recommendations)
analytics_b.add_edge(START, "analyze")
analytics_b.add_edge("analyze", "recommend")
analytics_b.add_edge("recommend", END)
analytics_sub = analytics_b.compile()
\end{lstlisting}
\end{frame}
%%%%%%%%%%%%%%%%%%%%%%%%%%%%%%%%%%%%%%%%%%%%%%%%%%%%%%%%%%%%%%%%%%%%%%%%%%%%%%%%
\begin{frame}[fragile]\frametitle{Practical Case: Marketing Agent, Main Workflow}

\begin{lstlisting}[language=Python]
def setup_campaign(state: CampaignState):
    # Fill in defaults for anything the caller left out
    return {"audience": state.get("audience") or "general",
            "platform": state.get("platform") or "twitter",
            "goal": state.get("goal") or "engagement"}

main_b = StateGraph(CampaignState)
main_b.add_node("setup", setup_campaign)
main_b.add_node("content", content_sub)          # subgraph as node
main_b.add_node("analytics", analytics_sub)      # subgraph as node

main_b.add_edge(START, "setup")
main_b.add_edge("setup", "content")
main_b.add_edge("content", "analytics")
main_b.add_edge("analytics", END)

marketing_agent = main_b.compile()
result = marketing_agent.invoke({"goal": "product launch"})
print(result["recommendations"])
\end{lstlisting}
\end{frame}
%%%%%%%%%%%%%%%%%%%%%%%%%%%%%%%%%%%%%%%%%%%%%%%%%%%%%%%%%%%%%%%%%%%%%%%%%%%%%%%%
\begin{frame}[fragile]\frametitle{Subgraph Best Practices and Considerations}

      \begin{itemize}
        \item Design Principles for Subgraphs
		      \begin{itemize}
				\item Keep subgraph functionality singular
				\item Ensure clear input and output interfaces
				\item Properly handle state passing
			  \end{itemize}

        \item Performance Considerations
		      \begin{itemize}
				\item Avoid frequent large data transfers between subgraphs
				\item Design state storage structures reasonably
				\item Consider asynchronous processing needs
			  \end{itemize}

        \item Error Handling
		      \begin{itemize}
				\item Implement error handling within subgraphs
				\item Provide clear error messages
				\item Ensure state consistency
			  \end{itemize}

      \end{itemize}
\end{frame}
%%%%%%%%%%%%%%%%%%%%%%%%%%%%%%%%%%%%%%%%%%%%%%%%%%%%%%%%%%%%%%%%%%%%%%%%%%%%%%%%
\begin{frame}[fragile]\frametitle{Multi-Agent Systems: Problems and Solutions}
      \begin{itemize}
          \item Agent is a system using LLM to control application flow
          \item Complex systems become harder to manage and scale over time
          \item Problem: Agents have too many tools, make poor decisions
          \item Problem: Context grows too complex for single agent tracking
          \item Problem: Need for multiple specialization areas (planning, research, math)
          \item Solution: Break application into smaller independent agents
          \item Independent agents can be simple prompt+LLM or complex React agents
          \item Primary benefit: Modularity makes development, testing, maintenance easier
          \item Primary benefit: Specialization creates expert domain-focused agents
          \item Primary benefit: Control explicit agent communication vs function calling
      \end{itemize}
\end{frame}
%%%%%%%%%%%%%%%%%%%%%%%%%%%%%%%%%%%%%%%%%%%%%%%%%%%%%%%%%%%%%%%%%%%%%%%%%%%%%%%%
\begin{frame}[fragile]\frametitle{Multi-Agent Architecture Types}
      \begin{itemize}
          \item Single Agent: Basic LLM with tools attachment
          \item Network: Four agents communicating with each other directly
          \item Network: Agents solve task parts, hand control to next agent
          \item Supervisor: One supervisor orchestrates team of specialized employees
          \item Supervisor: Coding agent, research agent, validation agent examples
          \item Supervisor: Agents don't communicate directly, only through supervisor
          \item Supervisor as Tools: Same as supervisor but agents provided as LLM tools
          \item Hierarchical: Extension of supervisor with multiple supervisor levels
          \item Hierarchical: Similar to big company organizational structure
          \item Custom: No clear definitions, conditional flow based on agent outputs
      \end{itemize}
\end{frame}
%%%%%%%%%%%%%%%%%%%%%%%%%%%%%%%%%%%%%%%%%%%%%%%%%%%%%%%%%%%%%%%%%%%%%%%%%%%%%%%%
\begin{frame}[fragile]\frametitle{Multi-Agent Architecture Types}

\begin{center}
\includegraphics[width=\linewidth,keepaspectratio]{langgraph12}

{\tiny (Ref: LangGraph Crash Course - Harish Neel)}

\end{center}


\end{frame}
%%%%%%%%%%%%%%%%%%%%%%%%%%%%%%%%%%%%%%%%%%%%%%%%%%%%%%%%%%%%%%%%%%%%%%%%%%%%%%%%
\begin{frame}[fragile]\frametitle{Supervisor Pattern in Code}
      \begin{itemize}
        \item The supervisor reads the shared state and picks the next worker, or finishes.
        \item \texttt{Command(goto=...)} performs the hand-off; workers always report back to the supervisor.
      \end{itemize}
% [code-font: forced \tiny | overflow at the default size]
% [code-font: blank lines removed from the listing to fit the frame]
\begin{lstlisting}[language=Python, basicstyle=\tiny\ttfamily]
from typing import Literal
from pydantic import BaseModel
from langgraph.graph import StateGraph, MessagesState, START, END
from langgraph.types import Command
class Route(BaseModel):
    next: Literal["researcher", "coder", "FINISH"]
def supervisor(state: MessagesState) -> Command[Literal["researcher", "coder", "__end__"]]:
    decision = llm.with_structured_output(Route).invoke(state["messages"])
    return Command(goto=END if decision.next == "FINISH" else decision.next)
def make_worker(agent):
    def worker(state: MessagesState) -> Command[Literal["supervisor"]]:
        result = agent.invoke(state)                # an inner agent
        return Command(update={"messages": result["messages"][-1:]},
                       goto="supervisor")           # always report back
    return worker
builder = StateGraph(MessagesState)
builder.add_node("supervisor", supervisor)
builder.add_node("researcher", make_worker(research_agent))
builder.add_node("coder", make_worker(code_agent))   # every goto target must exist
builder.add_edge(START, "supervisor")
graph = builder.compile()
\end{lstlisting}
\end{frame}
%%%%%%%%%%%%%%%%%%%%%%%%%%%%%%%%%%%%%%%%%%%%%%%%%%%%%%%%%%%%%%%%%%%%%%%%%%%%%%%%
\begin{frame}[fragile]\frametitle{Subgraphs: Foundation for Multi-Agent Systems}
      \begin{itemize}
          \item Subgraphs build complex systems with multiple graph components
          \item Common use case: Building multi-agent systems architecture
          \item Key question: How parent graph and subgraph communicate
          \item Key question: How to pass state between graphs during execution
          \item Scenario 1: Parent and subgraph share same schema keys
          \item Scenario 1: Add node with compiled subgraph directly
          \item Scenario 2: Parent and subgraph have different schemas
          \item Scenario 2: Add node function that invokes subgraph with transformation
          \item Transformation useful when state schemas differ between graphs
          \item Foundation enables building complex multi-agent architectures
      \end{itemize}
\end{frame}
%%%%%%%%%%%%%%%%%%%%%%%%%%%%%%%%%%%%%%%%%%%%%%%%%%%%%%%%%%%%%%%%%%%%%%%%%%%%%%%%
\begin{frame}[fragile]\frametitle{Implementing Subgraphs: Code Examples}
      \begin{itemize}
          \item Case 1: Shared schema - direct embedding approach
          \item Case 2: Different schemas - intermediate transformation node required
          \item Child graph: Simple LLM with Tavily search tool integration
          \item Parent graph: Embeds child graph as specialized search agent
          \item Shared schema: Compile subgraph, add as parent graph node
          \item Different schema: Extract query, transform to messages format
          \item Different schema: Invoke subgraph, extract response from result
          \item State transformation: Convert between parent and child schemas
          \item Example: Weather query routing through search agent subgraph
          \item Pattern enables modular, scalable multi-agent system construction
      \end{itemize}
\end{frame}
%%%%%%%%%%%%%%%%%%%%%%%%%%%%%%%%%%%%%%%%%%%%%%%%%%%%%%%%%%%%%%%%%%%%%%%%%%%%%%%%
\begin{frame}[fragile]\frametitle{Implementing Subgraphs: Shared Schema}

Parent and subgraph have shared state keys in their state schemas. In this case, you can include the subgraph as a node in the parent graph

% [code-font: forced \scriptsize | overflow at the default size]
\begin{lstlisting}[language=Python, basicstyle=\scriptsize\ttfamily]
from langgraph.graph import StateGraph, MessagesState, START

# Subgraph

def call_model(state: MessagesState):
    response = model.invoke(state["messages"])
    return {"messages": [response]}

subgraph_builder = StateGraph(MessagesState)
subgraph_builder.add_node("call_model", call_model)
subgraph_builder.add_edge(START, "call_model")
subgraph = subgraph_builder.compile()

# Parent graph

builder = StateGraph(MessagesState)
builder.add_node("subgraph_node", subgraph)
builder.add_edge(START, "subgraph_node")
graph = builder.compile()

graph.invoke({"messages": [{"role": "user", "content": "hi!"}]})
      \end{lstlisting}
\end{frame}
%%%%%%%%%%%%%%%%%%%%%%%%%%%%%%%%%%%%%%%%%%%%%%%%%%%%%%%%%%%%%%%%%%%%%%%%%%%%%%%%
\begin{frame}[fragile]\frametitle{Implementing Subgraphs: Different Schemas}

Parent graph and subgraph have different schemas (no shared state keys). In this case, you have to call the subgraph from inside a node in the parent graph

% [code-font: forced \tiny | overflow at the default size]
% [code-font: blank lines removed from the listing to fit the frame]
\begin{lstlisting}[language=Python, basicstyle=\tiny\ttfamily]
from typing_extensions import TypedDict, Annotated
from langchain_core.messages import AnyMessage
from langgraph.graph import StateGraph, MessagesState, START
from langgraph.graph.message import add_messages
class SubgraphMessagesState(TypedDict):
    subgraph_messages: Annotated[list[AnyMessage], add_messages]
# Subgraph
def call_model(state: SubgraphMessagesState):
    response = model.invoke(state["subgraph_messages"])
    return {"subgraph_messages": [response]}
subgraph_builder = StateGraph(SubgraphMessagesState)
subgraph_builder.add_node("call_model_from_subgraph", call_model)
subgraph_builder.add_edge(START, "call_model_from_subgraph")
subgraph = subgraph_builder.compile()
# Parent graph: a wrapper node maps the schemas
def call_subgraph(state: MessagesState):
    response = subgraph.invoke({"subgraph_messages": state["messages"]})
    return {"messages": response["subgraph_messages"]}
builder = StateGraph(MessagesState)
builder.add_node("subgraph_node", call_subgraph)
builder.add_edge(START, "subgraph_node")
graph = builder.compile()
graph.invoke({"messages": [{"role": "user", "content": "hi!"}]})
      \end{lstlisting}
\end{frame}
